% chapters/02_grundlagen.tex
%
% Grundlagenkapitel.
% Benötigt in der Präambel: amsmath, graphicx, booktabs.

\chapter{Grundlagen}
\label{chap:grundlagen}

Dieses Kapitel führt die Konzepte ein, die das Versuchsdesign in
Kapitel~\ref{chap:methodik} voraussetzt.
Es erklärt Begriffe und Mechanismen; die Datensätze sowie die
Einordnung einzelner Forschungsarbeiten und ihrer Ergebnisse folgen in
Kapitel~\ref{chap:verwandte_arbeiten}.

\section{Natural Language Processing (NLP)}
\label{sec:nlp}

Natural Language Processing (NLP) umfasst Verfahren, mit denen Computer
natürliche Sprache in geschriebener oder gesprochener Form verarbeiten,
analysieren und erzeugen~\cite{jurafsky2024speech,manning1999foundations}.
Typische Aufgaben sind die Klassifikation ganzer Texte, die Zuweisung von
Labels an einzelne Wörter einer Sequenz, die Extraktion strukturierter
Informationen und die Erzeugung von Text.
Named Entity Recognition gehört zur zweiten Gruppe und wird in
Abschnitt~\ref{sec:ner} eingeführt.

Heutige Systeme beruhen auf kontextualisierten Repräsentationen: Ein
neuronales Netz liest die gesamte Eingabesequenz und berechnet für jedes
Token einen Vektor, der von den umgebenden Tokens
abhängt~\cite{devlin2019bert}.
Daraus ist das heute dominierende Vorgehen in zwei Schritten
entstanden: Ein Modell wird zunächst auf großen, nicht annotierten
Textkorpora vortrainiert und anschließend an eine konkrete Aufgabe
angepasst, sei es durch Fine-Tuning seiner Gewichte oder durch
Anweisungen im Eingabetext~\cite{devlin2019bert,brown2020language}.
Vortrainierte Modelle lassen sich dabei auf zwei Arten nutzen.
Encoder liefern Repräsentationen, auf denen ein Klassifikator Labels
vorhersagt (Abschnitt~\ref{sec:encoder_ner}); generative Modelle erzeugen
die Lösung einer Aufgabe als Text (Abschnitt~\ref{sec:llms}).

\section{Named Entity Recognition (NER)}
\label{sec:ner}

Named Entity Recognition (NER) ist die Aufgabe, Erwähnungen benannter
Entitäten in einem Text zu lokalisieren und einer vordefinierten
Kategorie zuzuordnen~\cite{li2018survey}.
Formal erhält ein NER-System eine Tokenfolge $t_1, \dots, t_n$ und gibt
eine Menge von Tripeln $(I_s, I_e, \tau)$ zurück, wobei $I_s$ und $I_e$
den Start- und Endindex einer Erwähnung und $\tau$ ihren Typ
bezeichnen; die Formalisierung folgt \textcite{li2018survey}.
Eine Erwähnung ist damit ein zusammenhängender Abschnitt der Tokenfolge,
ein \emph{Span}; korrekt ist eine Vorhersage nur, wenn Grenzen und Typ
stimmen.
Das Prinzip zeigt Abbildung~\ref{fig:ner-grundprinzip} an einem
Beispiel.
Ob der Endindex das letzte Token einschließt oder auf das Token danach
zeigt, ist Konvention; die Abbildung verwendet einschließende Indizes,
die LLM-Pipeline in Abschnitt~\ref{sec:prompt-design} einen
ausschließenden Endindex.

\begin{figure}[htbp]
  \centering
  % figures/ner_grundprinzip.tex -- Grundprinzip der NER (eigene Darstellung,
% TikZ). Beispiel wie in Tabelle 2.1; Endindizes einschliessend.
\begin{tikzpicture}[
  font=\sffamily\small,
  tok/.style={draw=gray!60, fill=gray!8, rounded corners=2pt,
    minimum width=1.75cm, minimum height=0.8cm, inner sep=3pt},
  per/.style={tok, fill=orange!22}, loc/.style={tok, fill=blue!14},
  box/.style={draw=gray!60, fill=green!12, rounded corners=3pt,
    minimum width=5.2cm, minimum height=1.25cm, align=center,
    font=\sffamily\normalsize},
  outbox/.style={draw=gray!60, fill=white, rounded corners=3pt, inner sep=7pt},
  arr/.style={-{Latex[length=2.5mm]}, thick, gray!60!black},
  lbl/.style={font=\sffamily\footnotesize, text=gray!70!black},
]
  \node[lbl, anchor=west] at (-4.95, 1.25) {Eingabe: Tokenfolge $t_1, \dots, t_5$};
  \foreach \i/\w/\s in {1/Barack/per, 2/Obama/per, 3/visited/tok, 4/Berlin/loc, 5/./tok} {
    \node[\s] (t\i) at ({(\i-3)*2.05}, 0) {\w};
    \node[lbl] at ({(\i-3)*2.05}, 0.68) {$t_\i$};
  }
  \node[box] (ner) at (0, -2.05) {Named Entity Recognition\\(NER)};
  \draw[arr] (0, -0.42) -- (ner.north);
  \node[outbox, anchor=north] (res) at (0, -3.35) {%
    \begin{tabular}{@{}l@{\hspace{1.4em}}l@{}}
      $(t_1,\; t_2,\; \mathrm{PER})$ & \colorbox{orange!22}{\strut Barack Obama}\\[2pt]
      $(t_4,\; t_4,\; \mathrm{LOC})$ & \colorbox{blue!14}{\strut Berlin}
    \end{tabular}};
  \draw[arr] (ner.south) -- (res.north);
  \node[lbl, anchor=north] at ($(res.south)+(0,-0.12)$)
    {Ausgabe: Tripel (Startindex, Endindex, Typ), Endindex einschließend};
\end{tikzpicture}

  \caption[Grundprinzip der Named Entity Recognition]{Grundprinzip der Named Entity Recognition: Aus einer
  Tokenfolge werden Erwähnungen als Tripel aus Startindex, Endindex und
  Typ extrahiert.
  Eigene Darstellung in Anlehnung an
  \textcite[Abb.~1]{li2018survey}; der Endindex gehört zur Erwähnung.}
  \label{fig:ner-grundprinzip}
\end{figure}

Welche Kategorien unterschieden werden, legt der jeweilige Datensatz
fest.
Klassische Benchmarks verwenden wenige grobe Typen wie Person,
Organisation und Ort, feingranulare Datensätze wie MultiNERD
unterscheiden mehr als ein Dutzend Typen
(Abschnitt~\ref{sec:datasets_corpora}).

In der Praxis wird NER häufig als Sequenzlabeling formuliert: Jedes Token
erhält genau ein Label, das zugleich die Segmentierung in Spans und deren
Typ kodiert~\cite{li2018survey}.
Das Tagset $\{\mathrm{I}, \mathrm{O}, \mathrm{B}\}$ geht auf
\textcite{ramshaw1995text} zurück.
In der verbreiteten Variante IOB2, häufig kurz BIO-Schema genannt,
beginnt jede Entität mit einem B-Tag, I markiert ihre weiteren Tokens
und O alle Tokens außerhalb~\cite{tjong1999representing}.
Für NER wird der Typ an das Tag angehängt, etwa \texttt{B-PER} und
\texttt{I-PER}.
Die Kodierung setzt voraus, dass sich Entitäten weder überlappen noch
ineinander verschachtelt sind; CoNLL-2003 trifft genau diese
Annahme~\cite{sang2003conll}, und auch die in dieser Arbeit verwendeten
Datensätze enthalten nur flache Annotationen.
Ein Beispiel zeigt Tabelle~\ref{tab:bio-beispiel}.
Diese Arbeit verwendet IOB2 in der Encoder-Pipeline
(Abschnitt~\ref{sec:encoder-pipeline}) ebenso wie in der Auswertung
(Abschnitt~\ref{sec:evaluation}).

\begin{table}[htbp]
  \centering
  \small
  \renewcommand{\arraystretch}{1.3}
  \caption[IOB2-Kodierung eines Satzes]{IOB2-Kodierung eines Satzes.}
  \label{tab:bio-beispiel}
  \begin{tabular}{lccccc}
    \toprule
    \textbf{Token} & Barack & Obama & visited & Berlin & . \\
    \midrule
    \textbf{IOB2} & \texttt{B-PER} & \texttt{I-PER}
                  & \texttt{O} & \texttt{B-LOC} & \texttt{O} \\
    
    \bottomrule
  \end{tabular}
\end{table}

Lange Zeit waren Conditional Random Fields (CRF) für diese
Formulierung verbreitet; sie modellieren die bedingte
Wahrscheinlichkeit der gesamten Labelfolge bei gegebener Eingabe, sodass
Übergänge zwischen benachbarten Labels mitgelernt werden und unzulässige
Folgen unwahrscheinlich werden~\cite{lafferty2001crf}.
Neuronale Sequenzmodelle wie das BiLSTM-CRF ersetzten von Hand
entworfene Merkmale durch gelernte
Repräsentationen~\cite{lample2016neural,ma2016end}.

\section{Transformer-Architektur}
\label{sec:transformer}

Der Transformer von \textcite{vaswani2017attention} beruht
vollständig auf Attention: Jedes Token kann in einem Schritt auf jede
andere Position der Sequenz zugreifen.

Grundbaustein ist die \emph{Self-Attention}.
Aus der Repräsentation jedes Tokens werden durch gelernte lineare
Abbildungen drei Vektoren erzeugt: eine Query, ein Key und ein Value.
Das Skalarprodukt zwischen der Query eines Tokens und den Keys aller
Tokens misst deren Relevanz; die Softmax-Funktion normiert die
Relevanzwerte zu Gewichten, und die Ausgabe für jedes Token ist die
gewichtete Summe der
Value-Vektoren~\cite{vaswani2017attention}.
Transformer verwenden mehrere parallele \emph{Köpfe} mit jeweils
eigenen Projektionen.
Da Attention die Reihenfolge der Tokens nicht berücksichtigt, wird die
Position jedes Tokens eigens kodiert
(Abschnitt~\ref{sec:encoder_ner}).

Das Originalmodell besteht aus einem Encoder und einem Decoder.
Eine Transformer-Schicht des Encoders besteht aus zwei Teilschichten: der
Multi-Head-Self-Attention und einem positionsweisen Feed-Forward-Netz,
das jeden Tokenvektor unabhängig transformiert.
Der Decoder ergänzt jede Schicht um eine dritte Teilschicht, die auf die
Encoder-Ausgabe zugreift, und maskiert seine Self-Attention so, dass
Position $i$ nur auf Positionen bis $i$ zugreifen kann.
Diese \emph{kausale Maske} stellt sicher, dass die Vorhersage eines
Tokens nur von den bereits bekannten Tokens abhängt.
Den Aufbau zeigt Abbildung~\ref{fig:transformer-architektur}.

\begin{figure}[htbp]
  \centering
  % figures/vaswani_transformer.tex -- Transformer-Architektur (eigene
% Darstellung in Anlehnung an Vaswani et al. 2017, TikZ).
\begin{tikzpicture}[
  font=\sffamily\footnotesize,
  blk/.style={draw=gray!60, rounded corners=2pt, minimum width=3.5cm,
    minimum height=0.64cm, align=center, inner sep=2pt},
  attn/.style={blk, fill=orange!22}, ffn/.style={blk, fill=blue!12},
  norm/.style={blk, fill=green!16, minimum height=0.5cm},
  emb/.style={blk, fill=red!10}, lin/.style={blk, fill=violet!14},
  smx/.style={blk, fill=yellow!22, minimum height=0.5cm},
  stack/.style={draw=gray!50, rounded corners=4pt, fill=gray!6, inner sep=5pt},
  arr/.style={-{Latex[length=2mm]}, thick, black!75},
  res/.style={-{Latex[length=2mm]}, thick, black!75, rounded corners=2pt},
  pe/.style={circle, draw=gray!60, inner sep=1pt, minimum size=0.42cm, fill=white},
]
  % ---------------- Encoder (x = 0)
  \node[emb]  (ein) at (0, -0.15) {Eingabe-Embedding};
  \node[attn] (esa) at (0, 1.9)  {Multi-Head-\\Self-Attention};
  \node[norm] (en1) at (0, 2.75) {Add \& Norm};
  \node[ffn]  (eff) at (0, 3.75) {Feed-Forward-\\Netz};
  \node[norm] (en2) at (0, 4.6)  {Add \& Norm};
  \begin{scope}[on background layer]
    \node[stack, fit={(esa)(en2)(-2.05, 1.3)}] (enc) {};
  \end{scope}
  \node[font=\sffamily\small, anchor=east] at ($(enc.west)+(-0.25,0)$) {$N\times$};
  \node (einp) at (0, -1.3) {Eingaben};
  \node[pe] (eplus) at (0, 0.72) {$+$};
  \node[pe] (epos) at (-2.0, 0.72) {$\sim$};
  \node[align=center, font=\sffamily\scriptsize, anchor=east] at (-2.3, 0.72) {Positions-\\kodierung};
  \draw[arr] (einp) -- (ein);
  \draw[arr] (ein) -- (eplus); \draw[arr] (epos) -- (eplus); \draw[arr] (eplus) -- (esa);
  \draw[arr] (esa) -- (en1); \draw[arr] (en1) -- (eff); \draw[arr] (eff) -- (en2);
  \draw[res] (0, 1.3)  -- ++(-2.05, 0) |- (en1.west);
  \draw[res] (0, 3.2)  -- ++(-2.05, 0) |- (en2.west);
  % ---------------- Decoder (x = 6.6)
  \node[emb]  (aus) at (6.6, -0.15) {Ausgabe-Embedding};
  \node[attn] (dma) at (6.6, 1.9)  {Maskierte Multi-Head-\\Self-Attention};
  \node[norm] (dn1) at (6.6, 2.75) {Add \& Norm};
  \node[attn] (dca) at (6.6, 3.75) {Encoder-Decoder-\\Attention};
  \node[norm] (dn2) at (6.6, 4.6)  {Add \& Norm};
  \node[ffn]  (dff) at (6.6, 5.6)  {Feed-Forward-\\Netz};
  \node[norm] (dn3) at (6.6, 6.45) {Add \& Norm};
  \begin{scope}[on background layer]
    \node[stack, fit={(dma)(dn3)(8.65, 1.3)}] (dec) {};
  \end{scope}
  \node[font=\sffamily\small, anchor=west] at ($(dec.east)+(0.25,0)$) {$N\times$};
  \node[lin]  (lin)  at (6.6, 7.55) {Lineare Schicht};
  \node[smx]  (soft) at (6.6, 8.35) {Softmax};
  \node[align=center] (outp) at (6.6, 9.3) {Ausgabe-\\wahrscheinlichkeiten};
  \node[align=center] (ausp) at (6.6, -1.3) {Ausgaben\\(rechtsverschoben)};
  \node[pe] (dplus) at (6.6, 0.72) {$+$};
  \node[pe] (dpos) at (8.6, 0.72) {$\sim$};
  \node[align=center, font=\sffamily\scriptsize, anchor=west] at (8.9, 0.72) {Positions-\\kodierung};
  \draw[arr] (ausp) -- (aus);
  \draw[arr] (aus) -- (dplus); \draw[arr] (dpos) -- (dplus); \draw[arr] (dplus) -- (dma);
  \draw[arr] (dma) -- (dn1); \draw[arr] (dn1) -- (dca); \draw[arr] (dca) -- (dn2);
  \draw[arr] (dn2) -- (dff); \draw[arr] (dff) -- (dn3); \draw[arr] (dn3) -- (lin);
  \draw[arr] (lin) -- (soft); \draw[arr] (soft) -- (outp);
  \draw[res] (6.6, 1.3)  -- ++(2.05, 0) |- (dn1.east);
  \draw[res] (6.6, 3.2)  -- ++(2.05, 0) |- (dn2.east);
  \draw[res] (6.6, 5.05) -- ++(2.05, 0) |- (dn3.east);
  % ---------------- Encoder-Ausgabe -> Encoder-Decoder-Attention (Keys, Values)
  \draw[arr] (en2.north) -- ++(0, 0.45) -| (3.3, 3.75) -- (dca.west);
  \node[font=\sffamily\scriptsize, anchor=south, inner sep=2pt] at (1.65, 5.3) {Encoder-Ausgabe (Keys, Values)};
\end{tikzpicture}

  \caption[Transformer-Architektur mit Encoder und Decoder]{Transformer-Architektur mit Encoder (links) und Decoder
  (rechts).
  Jede Encoder-Schicht enthält Multi-Head-Self-Attention und ein
  Feed-Forward-Netz, jeweils mit Residualverbindung und
  Layer-Normalisierung; der Decoder ersetzt die Self-Attention durch eine
  maskierte und ergänzt eine Attention über die Encoder-Ausgabe.
  Eigene Darstellung in Anlehnung an \textcite[Abb.~1]{vaswani2017attention}.}
  \label{fig:transformer-architektur}
\end{figure}

Aus diesem Bauplan sind drei Modellfamilien hervorgegangen.
\emph{Encoder-only}-Modelle wie BERT verwenden nur den Encoder; jedes
Token sieht den Kontext in beide Richtungen~\cite{devlin2019bert}.
\emph{Decoder-only}-Modelle wie GPT-3 verwenden nur den Decoder ohne
Encoder-Attention; durch die kausale Maske eignen sie sich zur
Texterzeugung Token für Token~\cite{brown2020language}.
Encoder-Decoder-Modelle behalten den Aufbau des Originals bei.
Für NER folgt daraus der zentrale Unterschied der beiden verglichenen
Modellklassen: Ein Encoder liefert für jedes Token eine
beidseitig kontextualisierte Repräsentation, aus der ein Label
klassifiziert werden kann.
Ein Decoder verarbeitet die Sequenz kausal und gibt in der hier
verwendeten Pipeline die Entitäten als Text aus. Da der ganze
Eingabesatz im Prompt liegt, kann er dabei dennoch auch den Kontext
rechts von einer Entität nutzen.

Die Kosten der Self-Attention wachsen quadratisch mit der
Sequenzlänge~\cite{vaswani2017attention}.
Bei der Texterzeugung werden zudem die Keys und Values aller bisherigen
Tokens zwischengespeichert; Grouped-Query-Attention lässt mehrere
Query-Köpfe denselben Key- und Value-Kopf nutzen und verkleinert so
diesen Zwischenspeicher~\cite{ainslie2023gqa}.
Lineare Attention-Verfahren wie Gated DeltaNet ersetzen die
Relevanzmatrix durch einen Zustand fester Größe, sodass der Aufwand nur
linear mit der Sequenzlänge wächst~\cite{yang2025gated}.
Die in dieser Arbeit verwendeten LLMs wechseln lineare Schichten mit
einzelnen Schichten voller Attention ab~\cite{qwen2026card08b,qwen2026card4b,qwen2026card27b}; diese
hybride Attention ist für die Platzierung der Adapter relevant
(Abschnitt~\ref{sec:lora-konfiguration}).

\section{Encoder-basierte Modelle für NER}
\label{sec:encoder_ner}

Vortrainierte Transformer-Encoder sind seit BERT~\cite{devlin2019bert}
die übliche Architektur für überwachte NER.
Bevor ein Encoder einen Text verarbeiten kann, muss dieser in Einheiten
aus einem festen Vokabular zerlegt werden.
Da ein Vokabular aus ganzen Wörtern seltene Wörter nicht abdecken kann,
arbeiten alle hier behandelten Modelle mit \emph{Subwort-Einheiten}:
Häufige Wörter bleiben ganz, seltene werden in mehrere Teilstücke
zerlegt.
DeBERTa-v3 verwendet ein
SentencePiece-Vokabular~\cite{kudo2018sentencepiece} mit
128\,000 Einheiten~\cite{he2021debertav3}.
Für NER hat die Zerlegung eine praktische Folge: Die Labels liegen je
Wort vor, die Modelleingabe besteht aber aus Subwörtern.
\textcite{devlin2019bert} verwenden deshalb die Repräsentation des ersten
Subworts eines Wortes als Eingabe des Token-Klassifikators; in der
üblichen Umsetzung erhalten die übrigen Subwörter kein Label und gehen
nicht in den Verlust ein.
Die Encoder-Pipeline dieser Arbeit folgt diesem Vorgehen
(Abschnitt~\ref{sec:datenvorverarbeitung}).

\textcite{devlin2019bert} trainieren BERT auf nicht annotiertem Text vor allem mit
\emph{Masked Language Modeling} (MLM): Ein Teil der Tokens wird maskiert
und aus dem beidseitigen Kontext rekonstruiert.
Für eine Zielaufgabe wird das vortrainierte Modell um eine
aufgabenspezifische Ausgabeschicht ergänzt und vollständig feinjustiert,
das heißt, alle Encoder-Gewichte und die neue Schicht werden gemeinsam
mit dem Verlust der Zielaufgabe optimiert.
Für NER ist diese Ausgabeschicht eine lineare Klassifikationsschicht, die
aus der Repräsentation $h_i$ jedes Tokens Werte $W h_i + b$ für alle
Labels des Tagsets berechnet; eine Softmax-Funktion macht daraus eine
Verteilung, und bei der Vorhersage wird das Label mit dem höchsten Wert
gewählt.
Trainiert wird mit der Kreuzentropie zwischen dieser Verteilung und dem
Referenzlabel, gemittelt über alle Tokens mit Label.
Da die Positionskodierung von BERT auf 512 Positionen festgelegt
ist~\cite[Anh.~A.2]{devlin2019bert},
müssen längere Eingaben gekürzt oder geteilt werden; auch für
DeBERTa-v3 ist die maximale Eingabelänge auf 512 Positionen konfiguriert
(Parameter \texttt{max\_position\_embeddings} in den
Konfigurationsdateien der Revisionen in
Tabelle~\ref{tab:app-revisionen}),
sodass eine Längenbegrenzung ebenso für die Encoder-Pipeline dieser Arbeit
gilt (Abschnitt~\ref{sec:datenvorverarbeitung}).

DeBERTa verändert zwei Komponenten des
Encoders~\cite{he2020deberta}.
Erstens trennt die \emph{Disentangled Attention} Inhalt und Position:
Jedes Token wird durch einen Inhaltsvektor und einen Vektor für seine
relative Position zu den anderen Tokens repräsentiert, sodass die
Attention getrennt lernen kann, welche Inhalte und welche Abstände
zusammengehören.
Zweitens bezieht der \emph{Enhanced Mask Decoder} die absoluten
Positionen erst unmittelbar vor der Vorhersage der maskierten Tokens
ein.
DeBERTa-v3 behält diese Architektur bei und ändert das Vortraining in
zwei Punkten~\cite{he2021debertav3}.
Statt MLM verwendet es \emph{Replaced Token Detection} nach ELECTRA: Ein
kleines Generatormodell ersetzt einige Tokens durch plausible
Alternativen, und das eigentliche Modell lernt als Diskriminator, für
jedes Token zu entscheiden, ob es ersetzt wurde.
Das Trainingssignal entsteht so an allen Positionen statt nur an den
maskierten, was das Vortraining dateneffizienter
macht~\cite{clark2020electra}.
Zudem teilen Generator und Diskriminator beim \emph{Gradient-Disentangled
Embedding Sharing} ihre Token-Embeddings, ohne dass die Gradienten des
Diskriminators in die gemeinsamen Embeddings
zurückfließen~\cite{he2021debertav3}.

Die Base-Variante hat zwölf Schichten mit 86~Mio.
Parametern im Backbone, die Large-Variante 24 Schichten mit 304~Mio.,
jeweils zuzüglich der Embedding-Parameter des großen
Vokabulars~\cite{he2021debertav3card,he2021debertav3cardlarge}.

\section{Large Language Models (LLMs)}
\label{sec:llms}

Large Language Models (LLMs) sind Sprachmodelle mit typischerweise
Milliarden von Parametern, die auf sehr großen Textkorpora vortrainiert werden und als einzelnes
Modell ein breites Spektrum von Aufgaben bearbeiten~\cite{brown2020language,zhao2023survey}.
Die heute verbreiteten LLMs sind Decoder-only-Transformer
(Abschnitt~\ref{sec:transformer}).
Sie modellieren die Wahrscheinlichkeit einer Tokenfolge autoregressiv
als Produkt bedingter Wahrscheinlichkeiten,
\begin{equation}
  p(x_1, \dots, x_T) = \prod_{t=1}^{T} p(x_t \mid x_1, \dots, x_{t-1}),
  \label{eq:lm}
\end{equation}
und werden darauf trainiert, an jeder Position das nächste Token
vorherzusagen~\cite{brown2020language}.
Text wird erzeugt, indem das Modell wiederholt eine Verteilung über das
nächste Token berechnet, ein Token auswählt und es an die Eingabe
anhängt.
Jedes erzeugte Token erfordert einen eigenen Vorwärtsdurchlauf, der mit
wachsendem Kontext teurer wird (Abschnitt~\ref{sec:transformer}).
Die Ausgabelänge ist neben der Eingabelänge ein wesentlicher Faktor der
Inferenzzeit, was bei generativer NER die Wahl des Ausgabeformats zu
einer Effizienzfrage macht (Abschnitt~\ref{sec:llm_ner_related}).
Wie das Token ausgewählt wird, legt das \emph{Decoding} fest.
Greedy-Decoding wählt stets das wahrscheinlichste Token und ist als
Auswahlregel deterministisch; Gleitkommaarithmetik und Batching können
hardwareseitig Abweichungen verursachen.
Beim Sampling wird ein Token zufällig aus der Verteilung gezogen; eine
Temperatur skaliert die Verteilung, Top-$k$ beschränkt die Auswahl auf
die $k$ wahrscheinlichsten Tokens, und Top-$p$ (\emph{Nucleus Sampling})
auf die kleinste Menge von Tokens, deren Wahrscheinlichkeiten sich zu
mindestens $p$ summieren~\cite{holtzman2020curious}.
Gesampelte Ausgaben variieren zwischen Durchläufen, was bei der
Interpretation von Ergebnissen zu berücksichtigen ist
(Abschnitt~\ref{sec:llm-pipeline}).

Das Vortraining liefert ein Basismodell, das Text fortsetzt, aber
Anweisungen nicht zuverlässig befolgt.
\textcite{ouyang2022instructgpt} zeigen, wie ein \emph{Nachtraining}
(\emph{Post-Training}) dies ändert, unter anderem durch überwachtes
Fine-Tuning auf von Menschen geschriebenen Antworten
(\emph{Supervised Fine-Tuning}, SFT).
Nachtrainierte Modelle verarbeiten Eingaben in einem Dialogformat aus
Nachrichten mit Rollen wie System, Nutzer und Assistent, das ein
\emph{Chat-Template} mit speziellen Steuertokens in eine Tokenfolge
überführt~\cite{qwen2025qwen3,qwen2026card27b}.
Qwen-Modelle ab Qwen3 können vor der eigentlichen Antwort einen
abgegrenzten Block mit Zwischenüberlegungen erzeugen; dieser
\emph{Thinking-Modus} lässt sich ein- und
ausschalten~\cite{qwen2025qwen3}.

Nachtrainierte LLMs lösen Aufgaben, die im Eingabetext beschrieben
werden: Beim \emph{In-Context Learning}~\cite{brown2020language} erhält
das Modell eine Anweisung und optional einige
Demonstrationsbeispiele im Prompt und löst die Aufgabe ohne Änderung
seiner Parameter.
Enthält der Prompt keine Beispiele, spricht man von Zero-Shot, sonst von
Few-Shot.
Einige Fähigkeiten treten dabei erst ab einer gewissen Modellgröße auf;
\textcite{wei2022emergent} bezeichnen sie als emergent.

Viele LLMs werden nur mit frei herunterladbaren Gewichten und einer
Lizenz veröffentlicht (\emph{offene Gewichte}, \emph{open weights}).
Zu dieser Gruppe gehört die Qwen-Familie von Alibaba mit den Generationen
Qwen2.5, Qwen3 und Qwen3.5~\cite{qwen2024qwen25,qwen2025qwen3}; zur im
Februar 2026 vorgestellten Generation
Qwen3.5~\cite{qwen2026qwen35blog} gehören die hier verwendeten Modelle
Qwen3.5-27B, Qwen3.5-4B und Qwen3.5-0.8B, deren Modellkarten im
Schichtaufbau Feed-Forward-Netze und keine
Mixture-of-Experts-Schichten ausweisen~\cite{qwen2026card27b,qwen2026card4b,qwen2026card08b}
(Abschnitt~\ref{sec:llm-auswahl}).

Für NER ergibt sich aus der generativen Arbeitsweise eine andere
Formulierung als beim Sequenzlabeling.
Das Modell erhält die Aufgabenbeschreibung und den Text als Prompt und
erzeugt die Entitäten als Text, etwa als strukturierte Liste im
JSON-Format~\cite{dagdelen2024structuredIE}, je nach Verfahren mit
Angaben zur Position~\cite[Abschn.~3.2]{semin2026spanLabeling}.
Diese Ausgabe muss anschließend auf die Tokenfolge abgebildet werden,
damit sie mit den Referenzlabels verglichen werden kann.
Ausgabeformat und Nachverarbeitung werden damit Teil des Systems, und
die Ausgabe kann formal fehlerhaft sein, etwa durch ungültige Indizes
oder Typen außerhalb des Inventars~\cite{semin2026spanLabeling}.
Verfahren des \emph{Constrained Decoding} binden die Erzeugung an eine
formale Grammatik~\cite{geng2023grammar}; ob die Entitäten inhaltlich
stimmen und im Eingabetext vorkommen, bleibt dabei ungesichert.
Alternativ wird frei generiert und die Ausgabe nachträglich geprüft
(Abschnitte~\ref{sec:llm_ner_related} und~\ref{sec:prompt-design}).

\section{Parameter-Efficient Fine-Tuning (LoRA und QLoRA)}
\label{sec:lora}

Vollständiges Fine-Tuning passt alle Gewichte eines vortrainierten
Modells an.
Bei LLMs stößt das an Speichergrenzen, weil während des Trainings zu
jedem Parameter auch sein Gradient und die Zustände des Optimierers
gehalten werden müssen; Adam speichert je Parameter zwei zusätzliche Größen.
\emph{Parameter-Efficient Fine-Tuning} (PEFT) hält die vortrainierten
Gewichte eingefroren und trainiert nur eine kleine Zahl zusätzlicher
oder ausgewählter Parameter~\cite{han2024parameter}.

LoRA (\emph{Low-Rank Adaptation}) trainiert dazu eine
niedrigdimensionale Darstellung der
Gewichtsänderung~\cite{hu2021lora}.
Für eine vortrainierte Gewichtsmatrix $W_0 \in \mathbb{R}^{d \times k}$
wird die Änderung $\Delta W$ durch das Produkt zweier kleiner Matrizen
dargestellt, sodass sich für eine Eingabe $x$ die Ausgabe
\begin{equation}
  h = W_0 x + \Delta W x = W_0 x + \frac{\alpha}{r}\, B A x,
  \qquad B \in \mathbb{R}^{d \times r},\;
  A \in \mathbb{R}^{r \times k},\;
  r \ll \min(d, k)
  \label{eq:lora}
\end{equation}
ergibt.
$W_0$ bleibt eingefroren; trainiert werden nur $A$ und $B$.
Der Rang $r$ bestimmt die Kapazität des Adapters und die Zahl der
trainierbaren Parameter $r\,(d + k)$ anstelle von $d \cdot k$; die
Konstante $\alpha$ skaliert den Beitrag des
Adapters~\cite[Abschn.~4.1]{hu2021lora}.
In gängigen Implementierungen dient ein Dropout auf dem Adapterpfad als
Regularisierung.
$A$ wird zufällig initialisiert und $B$ mit null, sodass $\Delta W$ zu
Beginn des Trainings null ist und das Modell unverändert
startet~\cite{hu2021lora}.
Den Aufbau zeigt Abbildung~\ref{fig:lora-adapter}.
Nach dem Training kann $\Delta W$ zu $W_0$ addiert werden, sodass
gegenüber dem Ausgangsmodell keine zusätzliche Inferenzlatenz entsteht;
alternativ bleibt der Adapter getrennt und wird bei der Inferenz
mitberechnet.
\textcite{hu2021lora} adaptieren in ihren Experimenten die Query- und
Value-Projektionen der Attention; \textcite{dettmers2023qlora} setzen
Adapter auf alle linearen Schichten
(Abschnitt~\ref{sec:lora-konfiguration}).
Auf RoBERTa, DeBERTa, GPT-2 und GPT-3 erreicht LoRA bei einem Bruchteil
der trainierbaren Parameter eine Qualität, die der des vollständigen
Fine-Tunings entspricht oder sie übertrifft~\cite{hu2021lora}.

\begin{figure}[htbp]
  \centering
  % Abbildung 2.3: Aufbau eines LoRA-Adapters (TikZ)
\begin{tikzpicture}[
  >=Latex, font=\small,
  frozen/.style={draw, rounded corners=2pt, fill=gray!15, minimum width=2.6cm, minimum height=1.1cm, align=center},
  train/.style={draw, rounded corners=2pt, fill=blue!10, minimum width=1.6cm, minimum height=0.9cm, align=center},
  io/.style={draw, circle, minimum size=0.75cm, inner sep=0pt},
  sum/.style={draw, circle, minimum size=0.5cm, inner sep=0pt},
  arr/.style={->, thick}]
  \node[io] (x) at (0,0) {$x$};
  \node[frozen] (w) at (3.4,1.1) {$W_0 \in \mathbb{R}^{d \times k}$\\[1pt]{\footnotesize eingefroren}};
  \node[train] (a) at (2.5,-1.2) {$A \in \mathbb{R}^{r \times k}$};
  \node[train] (b) at (4.7,-1.2) {$B \in \mathbb{R}^{d \times r}$};
  \node[sum] (plus) at (6.6,0) {$+$};
  \node[io] (h) at (8.3,0) {$h$};
  \draw[arr] (x) -- (1.1,0) |- (w);
  \draw[arr] (x) -- (1.1,0) |- (a);
  \draw[arr] (a) -- (b);
  \draw[arr] (w) -| (plus);
  \draw[arr] (b) -| (plus);
  \draw[arr] (plus) -- (h);
  \node[font=\footnotesize, anchor=south] at (3.6,-0.55) {Rang $r \ll \min(d,k)$};
  \node[font=\footnotesize, anchor=north] at (3.6,-1.72) {trainierbarer Low-Rank-Pfad $\Delta W = \tfrac{\alpha}{r} BA$};
  \node[font=\footnotesize, anchor=south] at (8.3,0.65) {$h = W_0 x + \tfrac{\alpha}{r} B A x$};
\end{tikzpicture}

  \caption[Aufbau eines LoRA-Adapters]{Aufbau eines LoRA-Adapters.
  Die eingefrorene Gewichtsmatrix $W_0$ wird um den trainierbaren
  Low-Rank-Pfad $\Delta W = \tfrac{\alpha}{r} BA$ mit Rang $r$ ergänzt; beide Pfade
  verarbeiten dieselbe Eingabe $x$, und ihre Ausgaben werden addiert
  (Gleichung~\ref{eq:lora}).
  Eigene Darstellung in Anlehnung an \textcite[Abb.~1]{hu2021lora}.}
  \label{fig:lora-adapter}
\end{figure}

LoRA senkt den Speicherbedarf für Gradienten und Optimiererzustände,
nicht aber den für die eingefrorenen Gewichte selbst.
Hier setzt die \emph{Quantisierung} an.
Sie bildet Gewichte, die üblicherweise mit 16 oder 32~Bit gespeichert
werden, auf eine kleine Zahl von Stufen ab, etwa 256 Stufen bei 8~Bit.
Dazu wird eine Gruppe von Gewichten durch eine Skalierungskonstante auf
den Wertebereich des Zieltyps normiert; zur Berechnung wird sie
mit dieser Konstante wieder in höhere Präzision
zurückgerechnet~\cite{dettmers2022llmint8}.
Der Speicherbedarf sinkt entsprechend der Bitbreite.

QLoRA verbindet beides: Das Basismodell wird eingefroren und mit 4~Bit
gespeichert, gerechnet wird in 16~Bit, und die LoRA-Adapter werden in
16-Bit-Präzision trainiert~\cite{dettmers2023qlora}.
\textcite{dettmers2023qlora} führen dazu drei Komponenten ein.
Der Datentyp \emph{4-Bit-NormalFloat} (NF4) legt seine 16 Stufen so
fest, dass bei normalverteilten Gewichten in jede Stufe gleich viele
Werte fallen; er nutzt die vier Bit damit besser aus als
gleichabständige Stufen.
\emph{Double Quantization} quantisiert zusätzlich die
Skalierungskonstanten der Blöcke, die bei kleinen Blöcken selbst
spürbar Speicher belegen.
\emph{Paged Optimizers} lagern die Optimiererzustände bei
Speicherspitzen automatisch in den Hauptspeicher aus und verringern so
das Risiko von Abbrüchen wegen Speichermangels.
Im Vorwärts- und Rückwärtsdurchlauf werden die 4-Bit-Gewichte blockweise
in den Rechendatentyp BFloat16 zurückgerechnet; die Gradienten fließen
durch das eingefrorene quantisierte Modell, aktualisiert werden nur die
Adapter.
Mit diesem Verfahren lässt sich ein Modell mit 65~Mrd. Parametern auf
einer einzelnen GPU mit 48~GB feinjustieren; auf den von
\textcite{dettmers2023qlora} geprüften Benchmarks erreicht QLoRA die
Aufgabenleistung des 16-Bit-Fine-Tunings (volles Fine-Tuning bei GLUE
und Super-NaturalInstructions, LoRA bei MMLU).
In dieser Arbeit werden die Encoder vollständig feinjustiert und die
LLMs mit QLoRA angepasst; welche der genannten Komponenten dabei aktiv
sind, dokumentiert Abschnitt~\ref{sec:lora-konfiguration}.

\clearpage  % 2.7 beginnt auf neuer Seite
\section{Modellgröße und Parameteranzahl}
\label{sec:model_size}

Sprachmodelle unterscheiden sich um Größenordnungen in ihrer
Parameterzahl.
Wie sich die Größe auf die Qualität auswirkt, beschreiben
\emph{Skalierungsgesetze} in Form von Potenzgesetzen: Der Verlust fällt
als Potenzfunktion der Parameterzahl, der Datenmenge und des
Rechenaufwands~\cite{kaplan2020scaling}.
Die Gesetze beschreiben allerdings den Verlust des Vortrainings von
Modellen, die eigens für die Untersuchung trainiert wurden; über
Kennzahlen einer nachgelagerten Aufgabe wie den F$_1$-Score bei NER
sagen sie nichts unmittelbar aus.
Empirische Befunde zum Größeneinfluss bei NER fasst
Abschnitt~\ref{sec:model_size_related} zusammen; die Folgen für den
Größenvergleich dieser Arbeit nennt
Abschnitt~\ref{sec:einfluss-parameteranzahl}.

\section{Evaluationsmetriken}
\label{sec:evaluation}

NER-Systeme werden mit Precision, Recall und dem F-Maß bewertet, die
aus dem Information Retrieval
stammen~\cite{vanrijsbergen1979information,manning2008introduction}.
Auf Entitätsebene bezeichnet $TP$ die Zahl korrekt erkannter Entitäten,
$FP$ die Zahl vorhergesagter Entitäten ohne Entsprechung in der Referenz
und $FN$ die Zahl der Referenzentitäten, die das System nicht gefunden
hat.
Precision ist der Anteil korrekter Vorhersagen an allen Vorhersagen,
Recall der Anteil gefundener Entitäten an allen Referenzentitäten:
\begin{equation}
  \text{Precision} = \frac{TP}{TP + FP}, \qquad
  \text{Recall}    = \frac{TP}{TP + FN}, \qquad
  F_1 = 2 \cdot \frac{\text{Precision} \cdot \text{Recall}}%
                      {\text{Precision} + \text{Recall}} .
  \label{eq:prf1}
\end{equation}
Das F-Maß, das auf das Effektivitätsmaß $E$ bei van Rijsbergen
zurückgeht, fasst beide in einer Kennzahl
zusammen~\cite{vanrijsbergen1979information,manning2008introduction}; der F$_1$-Score ist das
harmonische Mittel beider Größen und die in der NER-Forschung
üblicherweise berichtete Kennzahl~\cite{sang2003conll,keraghel2024recent}.

Entscheidend ist, was als Instanz gezählt wird.
Bei einer Evaluation auf Tokenebene wird jedes Tokenlabel einzeln
bewertet; bei einer Evaluation auf Entitätsebene ist die Instanz die Entität, und eine
Vorhersage gilt nur dann als korrekt, wenn Spangrenzen und Typ exakt
mit der Referenz übereinstimmen~\cite{sang2003conll}.
Diese strenge Zählweise ist seit CoNLL-2003 üblich.
\textcite{segura2013semeval} definieren für die SemEval-2013-Aufgabe zur
Erkennung von Arzneimittelinteraktionen daneben weniger strenge Modi, die
nur die exakten Grenzen, eine teilweise Überlappung der Grenzen oder den
Typ bei überlappenden Grenzen als Treffer werten.
Solche Modi machen partielle Treffer sichtbar, erschweren aber den
Vergleich zwischen Arbeiten; diese Arbeit verwendet ausschließlich den
strikten Modus (Abschnitt~\ref{sec:performanzmetriken}).

In der Praxis werden Metriken auf Entitätsebene aus Tagfolgen berechnet;
wie dabei Spans aus den Tags gewonnen werden, hängt von der verwendeten
Bibliothek ab.
Die in dieser Arbeit verwendete Bibliothek \texttt{seqeval} bietet zwei
Modi~\cite{nakayama2018seqeval}.
Der Standardmodus ahmt das CoNLL-Auswertungsskript nach und toleriert formal unzulässige Tagfolgen; ein
\texttt{I-}Tag ohne vorangehenden Entitätsbeginn eröffnet dort dennoch
eine Entität.
Der strikte Modus wertet die Tagfolgen nach einem explizit angegebenen
Schema wie IOB2 aus, sodass unzulässige Folgen keine Entität bilden.

Bei mehreren Entitätstypen werden die Kennzahlen aggregiert.
Der Mikro-F$_1$-Score summiert $TP$, $FP$ und $FN$ über alle Typen und
berechnet daraus eine Kennzahl, sodass häufige Typen das Ergebnis
dominieren; der Makro-F$_1$-Score mittelt die je Typ berechneten
F$_1$-Scores und gewichtet alle Typen gleich~\cite{manning2008introduction}.
Bei stark unausgewogenen Typverteilungen wie in MultiNERD, wo drei Typen
den größten Teil der Entitäten stellen, spiegelt der Mikro-F$_1$-Score
vor allem die Leistung auf diesen Typen wider; typbezogene Werte
ergänzen ihn (Abschnitt~\ref{sec:performanzmetriken}).
Für seltene Typen mit geringem Support, also wenigen
Referenzentitäten, sind typbezogene Werte allerdings mit großer Unsicherheit
behaftet, da wenige Fehler den Wert stark verändern.

\textcite{hripcsak2005agreement} zeigen zudem, dass das F-Maß beim
Vergleich zweier Annotationen der positiven spezifischen Übereinstimmung
entspricht.
Ein F$_1$-Score misst damit die Übereinstimmung mit der Referenz und
nur mittelbar die Richtigkeit; bei einer automatisch erzeugten Referenz
wie der von MultiNERD (Abschnitt~\ref{sec:datasets_corpora}) ist diese
Unterscheidung wesentlich.
